\documentclass[10pt]{article}

% ---------- Document Identity (update these only) ----------
\newcommand{\DocStage}{}
\newcommand{\DocTitle}{Your Road to Tier One \textbf{SOF}}
\newcommand{\ProgramName}{182-Week Civilian-to-Selection Readiness Pipeline}
\newcommand{\ProgramSubtitle}{}

% ---------- Page ----------
\usepackage[legalpaper,left=0.5in,right=0.5in,top=0.6in,bottom=0.45in]{geometry}

% ---------- Typography ----------
\usepackage[T1]{fontenc}
\usepackage[utf8]{inputenc}
\usepackage{libertinus}
\usepackage{microtype}
\usepackage{parskip}
\usepackage{ragged2e}
\usepackage{textcomp}
\AtBeginDocument{\color{black}}
\AtBeginDocument{\pagecolor{white}}
\AtBeginDocument{\justifying}

% ---------- Landscape pages ----------
\usepackage{pdflscape}

% ---------- Color ----------
\usepackage[table]{xcolor}
\definecolor{StageBlue}{HTML}{1F4E79}
\definecolor{StageBlueDark}{HTML}{163A5A}
\definecolor{LightGray}{HTML}{F2F4F7}
\definecolor{MidGray}{HTML}{D9DEE5}
\definecolor{RuleGray}{HTML}{C7CED8}
\definecolor{HeaderInk}{HTML}{000000}
\definecolor{Ink}{HTML}{000000}

% ---------- Utility ----------
\usepackage{enumitem}
\sloppy
\setlength{\emergencystretch}{2em}

% ---------- Boxes / UI ----------
\usepackage[most]{tcolorbox}

\tcbset{
  charterTitle/.style={
    enhanced,
    colback=StageBlue!4,
    colframe=StageBlue,
    boxrule=1.25pt,
    arc=3pt,
    left=16pt,right=16pt,top=14pt,bottom=14pt,
    borderline north={3.2pt}{0pt}{StageBlue},
    borderline west={4pt}{0pt}{StageBlue},
    borderline south={1.25pt}{0pt}{StageBlue},
    drop shadow={black!25!white},
  },
  charterRule/.style={
    enhanced,
    colback=LightGray,
    colframe=RuleGray,
    boxrule=0.85pt,
    arc=3pt,
    left=12pt,right=12pt,top=10pt,bottom=10pt,
    borderline west={3pt}{0pt}{StageBlue},
  },
  charterBadge/.style={
    enhanced,
    colback=StageBlue,
    coltext=white,
    colframe=StageBlueDark,
    boxrule=0.6pt,
    arc=2pt,
    left=6pt,right=6pt,top=3pt,bottom=3pt,
  }
}
\newtcbox{\Badge}{nobeforeafter,charterBadge}

% ---------- Lists ----------
\setlist[itemize]{leftmargin=1.5em, itemsep=0.25em, topsep=0.25em}
\setlist[enumerate]{leftmargin=1.8em, itemsep=0.25em, topsep=0.25em}

% ---------- TOC ----------
\usepackage{tocloft}
\setcounter{tocdepth}{3}
\setlength{\cftbeforesecskip}{3pt}
\setlength{\cftbeforesubsecskip}{2pt}
\setlength{\cftbeforesubsubsecskip}{0pt}
\renewcommand{\cftsecfont}{\bfseries}
\renewcommand{\cftsecpagefont}{\bfseries}
\renewcommand{\cftsubsecfont}{\normalfont}
\renewcommand{\cftsubsecpagefont}{\normalfont}
\renewcommand{\cftsubsubsecfont}{\normalfont}
\renewcommand{\cftsubsubsecpagefont}{\normalfont}
\setlength{\cftsecindent}{0pt}
\setlength{\cftsubsecindent}{1.25em}
\setlength{\cftsubsubsecindent}{2.8em}
\setlength{\cftsecnumwidth}{2.1em}
\setlength{\cftsubsecnumwidth}{2.8em}
\setlength{\cftsubsubsecnumwidth}{3.6em}

% Remove the built-in TOC heading so only the custom heading is shown
\makeatletter
\renewcommand{\tableofcontents}{\@starttoc{toc}}
\makeatother

% ---------- Header/Footer: page number only ----------
\usepackage{fancyhdr}
\pagestyle{fancy}
\fancyhf{}
\renewcommand{\headrulewidth}{0pt}
\renewcommand{\footrulewidth}{0pt}
\fancyfoot[C]{\thepage}

% ---------- Section formatting ----------
\usepackage{titlesec}

\titleformat{\section}
  {\Large\bfseries\color{Ink}}
  {\thesection}{0.6em}{}
  [\vspace{0.25em}\textcolor{StageBlue}{\rule{\linewidth}{0.8pt}}]

\titleformat{\subsection}
  {\large\bfseries\color{Ink}}
  {\thesubsection}{0.6em}{}

\titleformat{\subsubsection}
  {\normalsize\bfseries\color{Ink}}
  {\thesubsubsection}{0.6em}{}

\titlespacing*{\section}{0pt}{0.95em}{0.35em}
\titlespacing*{\subsection}{0pt}{0.65em}{0.25em}
\titlespacing*{\subsubsection}{0pt}{0.55em}{0.2em}

% ---------- Table engine ----------
\usepackage{tabularray}
\UseTblrLibrary{booktabs}

% ---------- tabularray: GLOBAL table treatment ----------
\SetTblrInner{
  cells = {fg=black, bg=white, valign=t},
}

% ---------- Header helpers ----------
\newcommand{\Hdr}[1]{\shortstack[c]{\strut #1\strut}}

% ---------- Lines ----------
\newcommand{\TitleLine}{\textcolor{StageBlue}{\rule{0.98\linewidth}{0.95pt}}}
\newcommand{\SubTitleLine}{\textcolor{RuleGray}{\rule{0.98\linewidth}{0.65pt}}}

% ---------- Continued on next page ----------
\newcommand{\ContinuedOnNextPageInline}{%
  \par
  \vfill
  \noindent\textcolor{RuleGray}{\rule{\linewidth}{1pt}}\vspace{-2pt}
  \noindent\hfill\textit{Continued on next page}\par\vspace{-10pt}
  \noindent\textcolor{RuleGray}{\rule{\linewidth}{1pt}}
  \clearpage
}

% ---------- PDF hyperlinks ----------
\usepackage[hidelinks]{hyperref}
\hypersetup{
  colorlinks=false,
  pdfborder={0 0 0},
  linktoc=all,
  pdftitle={\DocTitle},
  pdfauthor={\ProgramName}
}

% =========================================================
\begin{document}

\section*{9.D.1 Coach Implementation Guide}
\phantomsection
\addcontentsline{toc}{section}{9.D.1 Coach Implementation Guide}

Operational bullets, binding posture. No new doctrine.

\subsection*{Pre-launch review sequence}
\phantomsection
\addcontentsline{toc}{subsection}{Pre-launch review sequence}

\begin{itemize}
\item Confirm the canonical binder and folder bundle are synchronized and labeled v-1.0; confirm filenames follow [StageID]-[DeliverableID]\_[Title]\_\_v-1.0.pdf.
\item Read and apply in this order before first credited session:
  \begin{itemize}
  \item 9.C QA Summary and Known Issues (deployment status + Must Fix Pre-Deploy).
  \item 8.D Issue Log and Corrective Principles (\textbf{ISS-08}-\#\#\# authoritative).
  \item 1.F Governance Rules and Validity Doctrine (\textbf{VALID}/\textbf{MODIFIED}/\textbf{INVALID}; \textbf{VALID}/\textbf{INVALID}; gate outcomes; stop posture; patch discipline).
  \item 2.E Data Recording and Test Logistics Conventions (required fields, missing-data doctrine, environment IDs, Evidence Ref posture).
  \item 2.C Protocol Cards used in Phase 00 and early Phase 01 windows (C\# + invalidation triggers).
  \item 3.D Integrated Phase Calendars and 3.E Adjustment Rules (protected windows, reslot discipline, no stacking).
  \item 4.A–4.G Phase 00 specification and exit posture (Phase 00 boundaries, Week 4/8/12/16 Deload/Test posture).
  \end{itemize}
\item Verify open High issues are addressed by approved handling posture before execution; do not proceed with silent workarounds:
  \begin{itemize}
  \item \textbf{ISS-08-001} pull-up feasibility timing risk.
  \item \textbf{ISS-08-002} and \textbf{ISS-08-003} underwater governance feasibility/admissibility.
  \item \textbf{ISS-08-004} long-ruck compounded dependency stack.
  \end{itemize}
\end{itemize}

\subsection*{Medical clearance posture and stop-rule enforcement}
\phantomsection
\addcontentsline{toc}{subsection}{Medical clearance posture and stop-rule enforcement}

\begin{itemize}
\item Enforce clearance before any training is credited:
  \begin{itemize}
  \item clearance documentation must be embedded in the binder corpus before the first session counts toward gates.
  \end{itemize}
\item Apply stop rules and \textbf{MEDICAL} \textbf{HOLD} posture per Stage 1.F and Stage 0.5:
  \begin{itemize}
  \item stop rules apply to both sessions and tests.
  \item red-flag symptoms override schedule and performance.
  \end{itemize}
\item Authority rule:
  \begin{itemize}
  \item trainee reports immediately (symptoms, pain, environment hazards, equipment failures),
  \item coach decides once informed,
  \item safety overrides schedule and performance and cannot be negotiated.
  \end{itemize}
\end{itemize}

\subsection*{Phase 00 onboarding and validity compliance}
\phantomsection
\addcontentsline{toc}{subsection}{Phase 00 onboarding and validity compliance}

\begin{itemize}
\item Start-from-zero posture:
  \begin{itemize}
  \item enforce minimal kit and conservative exposure posture while preserving the 6-session rhythm.
  \end{itemize}
\item Session validity discipline (non-negotiable):
  \begin{itemize}
  \item enforce required session elements in correct order:
    \begin{itemize}
    \item Safety self-check → Warm-up → Mobility → Main work → Cardio → Cooldown → Recovery notes/post-session notes.
    \end{itemize}
  \item apply \textbf{VALID}/\textbf{MODIFIED}/\textbf{INVALID} tagging with reason-code posture (do not invent reason codes).
  \item treat incomplete logs as inadmissible for gates; do not “credit” sessions with missing required fields.
  \end{itemize}
\item Phase 00 testing posture:
  \begin{itemize}
  \item tests occur only in Weeks 4/8/12/16 Deload/Test windows as specified in Stage 4.
  \item if a scheduled test cannot be performed \textbf{VALID}, mark \textbf{INVALID} and reslot per Stage 3.E.
  \end{itemize}
\end{itemize}

\subsection*{Using Stage 2 tests and Stage 3 deload/test cadence}
\phantomsection
\addcontentsline{toc}{subsection}{Using Stage 2 tests and Stage 3 deload/test cadence}

\begin{itemize}
\item Tests are executed only under Stage 2 protocol cards and only in Stage 3 protected windows.
\item Invalid test provides no gate credit:
  \begin{itemize}
  \item treat as “not yet tested” for that gate requirement,
  \item reslot to the next protected window per 3.E without stacking.
  \end{itemize}
\item Enforce required logging fields per Stage 2.E for every test event:
  \begin{itemize}
  \item Route \textbf{ID} / Pool \textbf{ID} / Machine \textbf{ID} as applicable,
  \item tool identity (timer/scale/tape; tool description or \textbf{ID}),
  \item Warm-up Completed flag,
  \item Evidence Ref posture where required,
  \item conditions notes and stop/abort fields.
  \end{itemize}
\item Enforce swim unit discipline:
  \begin{itemize}
  \item Pool \textbf{ID} and pool length unit must be recorded; missing unit/Pool \textbf{ID} makes evidence inadmissible (\textbf{ISS-08-005}).
  \end{itemize}
\end{itemize}

\subsection*{Applying Stage 3 adjustment rules}
\phantomsection
\addcontentsline{toc}{subsection}{Applying Stage 3 adjustment rules}

\begin{itemize}
\item Downshift-not-drop-frequency posture:
  \begin{itemize}
  \item preserve the 6-session rhythm,
  \item reduce risk via modality substitution, intensity downshift, and volume/density reduction,
  \item reslot tests rather than compressing or stacking make-ups.
  \end{itemize}
\item Gate outcome defaults:
  \begin{itemize}
  \item \textbf{HOLD} is default when required evidence is missing/invalid or dependencies are unmet.
  \item \textbf{REGRESS} and \textbf{MEDICAL} \textbf{HOLD} only when explicit triggers apply per Stage 1.F.
  \end{itemize}
\item Underwater governance unavailable:
  \begin{itemize}
  \item do not schedule or attempt; record “not scheduled; inadmissible without governance” and continue surface-only posture (\textbf{ISS-08-002}; \textbf{ISS-08-003}).
  \end{itemize}
\end{itemize}

\subsection*{Weekly QA/log review discipline}
\phantomsection
\addcontentsline{toc}{subsection}{Weekly QA/log review discipline}

\begin{itemize}
\item Binding cadence:
  \begin{itemize}
  \item daily logging QC (same-day completion, missing fields marked \textbf{MISSING}/\textbf{UNKNOWN}),
  \item weekly review (validity distribution + compliance rollups),
  \item mid-phase and end-phase gate reviews only in protected windows.
  \end{itemize}
\item Gate decision packet discipline (embedded into binder corpus):
  \begin{itemize}
  \item decision date,
  \item outcome label (\textbf{ADVANCE}/\textbf{HOLD}/\textbf{REGRESS}/\textbf{MEDICAL} \textbf{HOLD}),
  \item admissible evidence summary (ST-*/C\#/\textbf{MET-2B}-\#\#\# references),
  \item disqualifiers reviewed (pain/gait flags, stop-rule events, environment hazards),
  \item reslot targets if applicable.
  \end{itemize}
\end{itemize}

\subsection*{Environment and equipment safety controls}
\phantomsection
\addcontentsline{toc}{subsection}{Environment and equipment safety controls}

\begin{itemize}
\item Before any phase introduces a Tier 1 requirement:
  \begin{itemize}
  \item verify Stage 6 tier timing readiness is met and Stage 7 overview minimum capability is present and stabilized outside protected windows.
  \end{itemize}
\item Session-level safety gating:
  \begin{itemize}
  \item visibility controls are operationally mandatory when low-light/traffic-adjacent; otherwise substitute to daylight/indoor (7.19).
  \item unsafe environment conditions (ice glaze, lightning, air-quality warnings, heat index concern) trigger substitution/reslot per Stage 3.E.
  \end{itemize}
\item Substitution posture:
  \begin{itemize}
  \item preserve intent, reduce risk, document substitutions.
  \item enforce environment IDs and tool consistency posture for any substitution that could affect comparability.
  \end{itemize}
\end{itemize}

\subsection*{Change control and versioning}
\phantomsection
\addcontentsline{toc}{subsection}{Change control and versioning}

\begin{itemize}
\item No mid-phase changes to doctrine without:
  \begin{itemize}
  \item patch record entry (\textbf{PATCH}-\#\#\# as applicable),
  \item validation check,
  \item revalidation scope determination per Stage 1.F.
  \end{itemize}
\item Cosmetic-only corrections allowed only when meaning is unchanged; still logged to preserve audit trail.
\item Any unresolved conflict remains a Known Issue referenced by \textbf{ISS-08}-\#\#\#; do not “fix in the margins.”
\end{itemize}

\subsection*{File hygiene and where logs live}
\phantomsection
\addcontentsline{toc}{subsection}{File hygiene and where logs live}

\begin{itemize}
\item Binding rule:
  \begin{itemize}
  \item runtime logs may be collected on paper or digitally, but must be exported to PDF and embedded into the binder corpus for gate and audit use.
  \end{itemize}
\item Storage posture:
  \begin{itemize}
  \item store exported logs and gate packets in the stage folders consistent with the 9.B folder tree.
  \item archive prior versions only in /99\_Change\_Log\_and\_Archive/ with Versioning artifacts in /99\_Change\_Log\_and\_Archive/Versioning/.
  \end{itemize}
\end{itemize}

\newpage

\section*{9.D.2 Trainee Quick-Start Guide}
\phantomsection
\addcontentsline{toc}{section}{9.D.2 Trainee Quick-Start Guide}

Operational bullets, binding posture. No new doctrine.

\subsection*{Required reading order}
\phantomsection
\addcontentsline{toc}{subsection}{Required reading order}

Read and use these Stage 9 documents in this order:

\begin{itemize}
\item 9.A Executive Program Snapshot.
\item 9.B Artifact Index and Folder Structure.
\item 9.C QA Summary and Known Issues (including Must Fix Pre-Deploy list).
\item 9.D Implementation Guides (this document).
\item 9.E Versioning Change-Log and Freeze Window Framework.
\item 9.F Final Deployment Checklists.
\end{itemize}

\subsection*{Clearance posture and stop rules}
\phantomsection
\addcontentsline{toc}{subsection}{Clearance posture and stop rules}

\begin{itemize}
\item Do not start training without medical clearance on file with the coach.
\item Stop rules apply immediately:
  \begin{itemize}
  \item report pain that changes mechanics,
  \item report red-flag symptoms (chest pressure/pain, fainting/near-fainting, unusual shortness of breath, new neurologic symptoms),
  \item report panic/dissociation or inability to self-regulate in-session,
  \item do not self-override.
  \end{itemize}
\item Coach has decision authority once informed. \textbf{HOLD} or \textbf{MEDICAL} \textbf{HOLD} may be invoked.
\end{itemize}

\subsection*{Commitment expectations}
\phantomsection
\addcontentsline{toc}{subsection}{Commitment expectations}

\begin{itemize}
\item Sessions per week must match the operating constraints in the package (6 sessions/week).
\item Every session must include all required elements in correct order; incomplete sessions/logs may be tagged \textbf{MODIFIED}/\textbf{INVALID} and may not count for gates.
\item Cardio block is mandatory as part of each session structure; follow coach direction and written constraints; log honestly.
\item If any conflict exists across documents, follow coach direction and cite the relevant \textbf{ISS-08}-\#\#\# from 9.C; do not improvise.
\end{itemize}

\subsection*{What to execute first}
\phantomsection
\addcontentsline{toc}{subsection}{What to execute first}

\begin{itemize}
\item Start with Phase 00 as packaged in Stage 4 (4.A–4.G) and follow the protected Week 4/8/12/16 Deload/Test posture.
\item Execute only what the phase prescribes at the intent level and only what your current equipment and environment can support safely.
\item Use logging conventions from Stage 2.E and validity tagging from Stage 1.F:
  \begin{itemize}
  \item record same-day,
  \item do not backfill by memory,
  \item missing fields remain \textbf{MISSING}/\textbf{UNKNOWN}.
  \end{itemize}
\end{itemize}

\subsection*{Gear posture}
\phantomsection
\addcontentsline{toc}{subsection}{Gear posture}

\begin{itemize}
\item Begin with the Phase 00 minimal kit posture and stabilize it; do not change shoes/tools inside protected windows.
\item Build capabilities only as Stage 6 indicates by phase and only after you can safely use them.
\item Do not attempt tasks requiring Tier 1 capabilities you do not have:
  \begin{itemize}
  \item report the gap immediately to the coach,
  \item do not improvise unsafe substitutes (example: ruck load carriage without stable ruck system; pull-ups without safe anchored setup).
  \end{itemize}
\end{itemize}

\subsection*{Sleep, hydration, and nutrition posture}
\phantomsection
\addcontentsline{toc}{subsection}{Sleep, hydration, and nutrition posture}

\begin{itemize}
\item Do not self-modify nutrition strategy as a program substitute; Stage 10 governs nutrition strategy once issued.
\item Treat sleep and hydration as execution enablers and log honestly:
  \begin{itemize}
  \item sleep hours,
  \item soreness/fatigue,
  \item pain flags and foot hotspots,
  \item steps (baseline activity) where required by the daily log.
  \end{itemize}
\end{itemize}

\subsection*{Reporting protocol}
\phantomsection
\addcontentsline{toc}{subsection}{Reporting protocol}

\begin{itemize}
\item Report immediately:
  \begin{itemize}
  \item pain that alters mechanics,
  \item illness symptoms,
  \item safety issues and near-misses,
  \item environment hazards (ice, lightning, air-quality alerts, unsafe heat),
  \item equipment failures (lights, ruck buckles, leaking reservoir).
  \end{itemize}
\item Coach decides once informed. \textbf{HOLD}, \textbf{REGRESS}, or \textbf{MEDICAL} \textbf{HOLD} may be invoked when indicated.
\item Do not hide deviations. Deviations must be logged so decisions remain auditable and gate outcomes remain defensible.
\end{itemize}

\end{document}